% Corrected continuation: static response and a positive photon-mass deformation.
\section{Static susceptibility and a positive photon-mass deformation}
\label{sec:ward}

This section gives two statements with different scopes. First, the free cutoff
Dirac sea has a negative second-order response to a sufficiently soft static
transverse field. Second, adding a fixed positive transverse photon-mass term
produces a modified Hamiltonian with an extensive lower bound. The second
statement does not establish stability of renormalized massless QED. The
distinction is necessary because the field-independent momentum cutoff used
here is not gauge covariant.

Write
\[
 E(p)=\sqrt{p^2+m^2},\qquad
 D(p)=\boldsymbol{\mathsf a}\cdot p+m\mathsf b,\qquad
 P_\pm(p)=\frac12\bigl(1\pm D(p)/E(p)\bigr),\qquad
 E_\Lambda=\sqrt{\Lambda_f^2+m^2}.
\]
All fermionic perturbation statements below concern \(H_D\) and its free
vacuum, without the Coulomb interaction. They are finite-dimensional
statements before any limit is taken.

\subsection{Static transverse susceptibility of the cutoff Dirac sea}

For \(q\in(2\pi/L)\Z^3\setminus\{0\}\) and a real unit vector
\(\varepsilon\perp q\), define
\begin{equation}
 \chi_L(q;\varepsilon)=\frac1V
 \sum_{\substack{p\in\Gamma_{f,L}\\p+q\in\Gamma_{f,L}}}
 \frac{\tr\left[P_+(p+q)(\varepsilon\cdot\boldsymbol{\mathsf a})
 P_-(p)(\varepsilon\cdot\boldsymbol{\mathsf a})\right]}
 {E(p+q)+E(p)}.
 \label{eq:chiL}
\end{equation}
Each summand is nonnegative. Indeed, with
\(X=P_+(p+q)(\varepsilon\cdot\boldsymbol{\mathsf a})P_-(p)\), its
numerator is \(\tr(XX^*)\).

\begin{lemma}[Trace identity]
\label{lem:trace}
For every \(p\in\R^3\) and every real unit vector \(\varepsilon\),
\begin{equation}
 \tr\left[P_+(p)(\varepsilon\cdot\boldsymbol{\mathsf a})
 P_-(p)(\varepsilon\cdot\boldsymbol{\mathsf a})\right]
 =2\left(1-\frac{(\varepsilon\cdot p)^2}{E(p)^2}\right).
 \label{eq:trace-identity}
\end{equation}
\end{lemma}
\begin{proof}
Set \(\mathsf e=\varepsilon\cdot\boldsymbol{\mathsf a}\). The Dirac
anticommutation relations give \(\mathsf e^2=1\) and
\[
 \mathsf eD(p)\mathsf e=-D(p)+2(\varepsilon\cdot p)\mathsf e.
\]
Using \(\tr D(p)=\tr\mathsf e=0\), \(D(p)^2=E(p)^2I\), and
\(\tr[D(p)\mathsf e]=4\varepsilon\cdot p\), we obtain
\[
 \tr(P_+\mathsf eP_-\mathsf e)
 =\frac14\left(4-\frac{\tr[D(p)\mathsf eD(p)\mathsf e]}{E(p)^2}\right)
 =\frac14\left(8-\frac{8(\varepsilon\cdot p)^2}{E(p)^2}\right).
\]
\end{proof}

\begin{proposition}[Second-order response to a static transverse field]
\label{prop:response}
Let \(A_{\rm cl}(x)=a\varepsilon\cos(q\cdot x)\), where \(a\in\R\), and put
\[
 J_q=\int_{\T_L^3}\cos(q\cdot x)\,\varepsilon\cdot j(x)\dd x,
 \qquad h(a)=H_D+eaJ_q.
\]
Restrict \(h(a)\) to the neutral fermionic sector. Then:
\begin{enumerate}[label=(\roman*)]
\item The free vacuum \(\Omega_f\) is the unique ground state of \(h(0)\),
and the remaining spectrum is contained in \([2m,\infty)\).
\item The lowest eigenvalue \(e_0(a)\) is real-analytic and even near zero,
and
\begin{equation}
 e_0(a)=-\frac{e^2a^2V}{2}\chi_L(q;\varepsilon)+O(a^4).
 \label{eq:second-order}
\end{equation}
The remainder in this statement is at fixed \(L,q\) and cutoffs.
\item For fixed \(n\in\Z^3\setminus\{0\}\), fixed \(\varepsilon\perp n\),
and \(q=2\pi n/L\),
\begin{equation}
 \lim_{L\to\infty}\chi_L(q;\varepsilon)=\chi(0)
 :=\int_{|p|\leq\Lambda_f}
 \frac{1-(\varepsilon\cdot p)^2/E(p)^2}{E(p)}\frac{\dd p}{(2\pi)^3}
 =\frac{\Lambda_f^3}{6\pi^2E_\Lambda}.
 \label{eq:chi0}
\end{equation}
\end{enumerate}
\end{proposition}
\begin{proof}
Every nonvacuum state in the neutral free sector contains at least one
electron and one positron, so its energy is at least \(2m\). The simplicity
and gap imply analytic perturbation theory for the lowest eigenvalue.
The vacuum current vanishes: its local sea contribution is an odd momentum
sum for the symmetric cutoff, and the nonzero cosine mode also has zero
momentum diagonal. Hence \(\langle\Omega_f,J_q\Omega_f\rangle=0\).

In the finite sea representation let
\[
 A_f=P_f(\varepsilon\cdot\boldsymbol{\mathsf a})
             \cos(q\cdot x)P_f.
\]
The normal-ordered bilinear \(J_q\) creates from the vacuum precisely the
particle-hole states with matrix elements
\(\langle\varphi_i,A_f\varphi_j\rangle\), where \(i\) is a positive-energy
state and \(j\) a negative-energy state of \(D_0\). If \(R_D\) denotes the
inverse of \(H_D\) on \(\Omega_f^\perp\) in the neutral sector, the
coefficient of \(a^2\) is
\[
 -e^2\langle J_q\Omega_f,R_DJ_q\Omega_f\rangle
 =-e^2\sum_{i,j}
   \frac{|\langle\varphi_i,A_f\varphi_j\rangle|^2}{E_i+|E_j|}.
\]
The two momentum transfers \(q\) and \(-q\) have matrix coefficient
\(\frac12\varepsilon\cdot\boldsymbol{\mathsf a}\). They give disjoint
particle-hole matrix entries because \(q\ne0\). Thus the sum is
\(V[\chi_L(q;\varepsilon)+\chi_L(-q;\varepsilon)]/4\).
The substitution \(p\mapsto-p\), together with
\[
 \mathsf bP_\pm(p)\mathsf b=P_\pm(-p),\qquad
 \mathsf b(\varepsilon\cdot\boldsymbol{\mathsf a})\mathsf b
 =-\varepsilon\cdot\boldsymbol{\mathsf a},
\]
proves \(\chi_L(-q;\varepsilon)=\chi_L(q;\varepsilon)\).
Finally, translation by any \(t\) with \(q\cdot t=\pi\) modulo \(2\pi\)
preserves the cutoff and \(H_D\), while sending \(J_q\) to \(-J_q\).
Consequently \(e_0(a)=e_0(-a)\), which proves the even remainder in
\eqref{eq:second-order}.

For (iii), the summand is bounded and continuous in its momenta, since
\(m>0\). The intersection of the two cutoff balls converges to the cutoff
ball, whose boundary has measure zero. Riemann-sum convergence and
Lemma~\ref{lem:trace} therefore give the integral in \eqref{eq:chi0}.
Angular averaging gives
\[
 \chi(0)=\frac1{2\pi^2}\int_0^{\Lambda_f}
 \left(\frac{p^2}{\sqrt{p^2+m^2}}
       -\frac{p^4}{3(p^2+m^2)^{3/2}}\right)\dd p
 =\frac{\Lambda_f^3}{6\pi^2E_\Lambda},
\]
as follows by differentiating \(p^3/(3\sqrt{p^2+m^2})\).
\end{proof}

\begin{corollary}[Negative static mass-like curvature]
\label{cor:tachyon}
The free-sea Born--Oppenheimer energy, including the classical magnetic
energy of the specified cosine mode, satisfies
\begin{equation}
 e_0(a)+\frac{V|q|^2a^2}{4}
 =\frac{Va^2}{4}\left(|q|^2-2e^2\chi_L(q;\varepsilon)\right)+O(a^4).
 \label{eq:tachyon}
\end{equation}
Its quadratic coefficient is negative exactly when
\[
 |q|^2<2e^2\chi_L(q;\varepsilon).
\]
Define the positive scale
\begin{equation}
 \mu_{\rm ind}^2:=2e^2\chi(0)
 =\frac{4\alpha}{3\pi}\frac{\Lambda_f^3}{E_\Lambda}
 =\frac{4\alpha}{3\pi}\Lambda_f^2
   \left(1+O(m^2/\Lambda_f^2)\right).
 \label{eq:muind}
\end{equation}
For fixed nonzero \(n\), fixed \(\varepsilon\perp n\), and any
\(0<\delta<1\), sufficiently large \(L\) satisfies
\(2e^2\chi_L(2\pi n/L;\varepsilon)>(1-\delta)\mu_{\rm ind}^2\).
In particular, the quadratic coefficient is negative along these sequences
once \(|2\pi n/L|^2<(1-\delta)\mu_{\rm ind}^2\).
\end{corollary}
\begin{proof}
For a transverse cosine, \(\frac12\int|\nabla\times A_{\rm cl}|^2
=V|q|^2a^2/4\). Apply Proposition~\ref{prop:response} and its limit.
\end{proof}

The limit used in this corollary does not imply instability for every
finite momentum with \(|q|<\mu_{\rm ind}\). The finite-momentum condition
contains \(\chi_L(q;\varepsilon)\), not \(\chi(0)\). The word ``mass-like''
describes the nonzero constant in this static quadratic coefficient. No
photon pole, real-time dispersion relation, or physical tachyon mass has
been derived. Coulomb and photon fluctuation corrections are not included
in this free-sea susceptibility.

\begin{remark}[Gauge interpretation and the scope of positivity]
\label{rem:ward}
On a torus a constant vector potential generally carries nontrivial
holonomy. It can be removed by a periodic charged gauge transformation only
if \(eA_iL\in2\pi\Z\) for every component. Thus finite-volume gauge
invariance does not by itself require independence of all constant
potentials. The model also fixes the harmonic photon sector by omitting
\(k=0\).

In an infinite-volume gauge-compatible treatment, the Ward identities,
together with the appropriate locality and regularity of the polarization
tensor near zero momentum, constrain an unphysical photon-mass term.
Those ingredients are not supplied by the present field-independent
momentum cutoff. In particular, the fact that the canonical Dirac coupling
is linear in \(A\) does not, by itself, disprove the possibility of such
cancellations after a different regularization and renormalization.
We do not assume an unproved Schwinger-term or diamagnetic cancellation.
The classical-field Pauli--Villars construction of \cite{GHLS2013} is an
example of a different framework; it is not a theorem for the present
quantized many-body Hamiltonian.

There is a valid abstract positivity statement. If a finite-dimensional
\(H_0\geq0\) has simple ground state \(\Omega\), a positive gap, and a
self-adjoint centered current \(J\), then
\[
 \left.\frac{\dd^2}{\dd a^2}\inf\sigma(H_0+eaJ)\right|_{a=0}
 =-2e^2\langle J\Omega,(H_0|_{\Omega^\perp})^{-1}J\Omega\rangle.
\]
This is strictly negative precisely when \(J\Omega\ne0\). In that case,
\[
 \langle J\Omega,(H_0|_{\Omega^\perp})^{-1}J\Omega\rangle
 \geq\frac{\|J\Omega\|^2}{\|H_0\|}>0.
\]
The hypotheses do not establish nonzero pair coupling or a positive
susceptibility per volume for arbitrary regulators. Here these facts are
verified by the explicit current and \eqref{eq:chi0}. The negative curvature
and the \(L^5\) trial estimate of Section~\ref{sec:counterexamples} both
exhibit the inadequacy of this massless, field-independent cutoff for a
volume-extensive vacuum energy.
\end{remark}

\paragraph{Numerical scale, not a finite-volume threshold theorem.}
For \(\alpha=1/137\), direct evaluation of \eqref{eq:chi0} and
\eqref{eq:muind} gives the following dimensionless values:
\begin{center}
\begin{tabular}{@{}cccc@{}}
\toprule
\(\Lambda_f/m\) & \(\chi(0)/m^2\) & \(\mu_{\rm ind}^2/m^2\)
& \((2\pi/\mu_{\rm ind})m\)\\
\midrule
5   & 0.413973 & 0.0759437 & 22.8000\\
10  & 1.680306 & 0.3082532 & 11.3169\\
100 & 168.860197 & 30.9775155 & 1.1289\\
\bottomrule
\end{tabular}
\end{center}
For \(\Lambda_f\gg m\), the continuum long-wavelength estimate for the
lowest-mode length scale is
\[
 L_*:=\frac{2\pi}{\mu_{\rm ind}}
 \sim\frac{112.8875}{\Lambda_f}.
\]
This estimates the scale of onset; an actual finite-box test must use
\eqref{eq:chiL} and the exact criterion in Corollary~\ref{cor:tachyon}.
No monotonicity in \(|q|\), or exact threshold inferred from a numerical
lattice sample, is asserted.

\subsection{An extensive bound for a positive photon-mass deformation}

Introduce an additional fixed parameter \(\delta\mu>0\) and define
\begin{equation}
 H^{(\delta\mu)}_{\Lambda,L}(R,Z)
 =H_{\Lambda,L}(R,Z)
 +\frac{\delta\mu^2}{2}\int_{\T_L^3}:\!A(x)\cdot A(x)\!:\dd x.
 \label{eq:Hct}
\end{equation}
Photon normal ordering here is relative to the original massless Fock
vacuum. The added operator is independent of states and nuclear
configurations. It changes the bare Hamiltonian. It is a transverse
photon-mass deformation of the two-polarization Coulomb-gauge model; the
instantaneous Coulomb interaction is unchanged. We do not identify it
with a complete covariant Proca theory, or assume that it is an allowed
renormalization of physical massless QED without further matching conditions.

Choose the real polarization frame so that
\(\varepsilon_\lambda(-k)=\varepsilon_\lambda(k)\). This is possible by
choosing the frame independently on each pair \(\{k,-k\}\); no smooth
global polarization frame is needed for these finite mode sets.

\begin{lemma}[Diagonalization of the modified photon Hamiltonian]
\label{lem:proca}
Let \(\omega_k=\sqrt{|k|^2+\delta\mu^2}\). There is a unitary \(U\) on the
photon Fock space, a finite product of oscillator squeezing transformations,
such that
\begin{align}
 U\left(H_\gamma+\frac{\delta\mu^2}{2}\int:\!A^2\!:\right)U^*
 &=\sum_{k,\lambda}\omega_k a^\dagger_{k,\lambda}a_{k,\lambda}
   +E_\gamma(\delta\mu),
 \label{eq:proca-H}\\
 U A(x)U^*
 &=\sum_{k,\lambda}\frac{\varepsilon_\lambda(k)}{\sqrt{2V\omega_k}}
 \left(a_{k,\lambda}e^{ik\cdot x}+a^\dagger_{k,\lambda}e^{-ik\cdot x}\right),
 \label{eq:proca-A}
\end{align}
where all sums run over the original retained nonzero photon modes and
\begin{equation}
 E_\gamma(\delta\mu)
 =\sum_{k,\lambda}
 \left(\frac{\omega_k-|k|}{2}-\frac{\delta\mu^2}{4|k|}\right)
 \in\left[-\frac{\delta\mu^2}{4}\sum_{k,\lambda}|k|^{-1},\,0\right].
 \label{eq:proca-const}
\end{equation}
For \(L\geq\Lambda_\gamma^{-1}\),
\( |E_\gamma(\delta\mu)|\leq C\delta\mu^2\Lambda_\gamma^2V\).
\end{lemma}
\begin{proof}
Use orthonormal real sine and cosine modes for each pair \(\{k,-k\}\)
and polarization. Each real oscillator has coordinates \(X,\Pi\),
\([X,\Pi]=i\), with massless Hamiltonian
\[
 \frac12(\Pi^2+|k|^2X^2)-\frac{|k|}{2},
 \qquad X=\frac{a+a^\dagger}{\sqrt{2|k|}}.
\]
The added normal-ordered term for that oscillator is
\(\frac{\delta\mu^2}{2}(X^2-1/(2|k|))\). Hence their sum is
\[
 \frac12(\Pi^2+\omega_k^2X^2)
 -\frac{|k|}{2}-\frac{\delta\mu^2}{4|k|}.
\]
The unitary dilation with
\[
 UXU^*=\sqrt{|k|/\omega_k}\,X,
 \qquad U\Pi U^*=\sqrt{\omega_k/|k|}\,\Pi
\]
sends the quadratic oscillator to
\(\omega_k(a^\dagger a+1/2)\), and sends its field coordinate to
\((a+a^\dagger)/\sqrt{2\omega_k}\). There are two real oscillators per
pair and polarization, exactly the number of original momentum-mode
oscillators. Returning to plane waves proves both operator identities and
the stated scalar constant, with no additional factor of two.

The identity
\(\omega_k-|k|=\delta\mu^2/(\omega_k+|k|)\) implies
\(0\leq\omega_k-|k|\leq\delta\mu^2/(2|k|)\), proving the scalar bounds.
Finally, lattice-shell counting gives
\[
 \frac1V\sum_{0<|k|\leq\Lambda_\gamma,\lambda}|k|^{-1}
 \leq C\Lambda_\gamma^2
 \qquad (L\geq\Lambda_\gamma^{-1}).
\]
For example, after writing \(k=2\pi n/L\), the sum of \(|n|^{-1}\)
over \(0<|n|\leq R\) is bounded by \(CR^2\); if \(R<1\), it is empty.
\end{proof}

\begin{theorem}[Extensive lower bound for the modified massive-bare model]
\label{thm:ct}
Let \(\delta\mu>0\) and
\(L\geq\max\{\Lambda_f^{-1},\Lambda_\gamma^{-1}\}\).
For every distinct nuclear configuration, the Hamiltonian
\eqref{eq:Hct} is self-adjoint on
\[
 (I\otimes U^*)D\left(I\otimes\sum_{k,\lambda}
                         \omega_k a^\dagger_{k,\lambda}a_{k,\lambda}\right)
\]
and has compact resolvent. The same statements hold on each nonempty
charge sector. On every admissible neutral sector,
\begin{equation}
 H^{(\delta\mu)}_{\Lambda,L}(R,Z)
 \geq-A_\Lambda K-
 \left(B_\Lambda+\frac{6e^2d_L^2}{\delta\mu^2}
                  +C\delta\mu^2\Lambda_\gamma^2\right)V,
 \label{eq:ct-bound}
\end{equation}
where \(A_\Lambda,B_\Lambda\) are the constants of
Section~\ref{sec:lower-proof} and \(d_L=M_L/V\).
Since \(d_L\leq C\Lambda_f^3\) in the specified range of box sizes,
this is a lower bound uniform in volume. Moreover,
\begin{equation}
 -\left(B_\Lambda+\frac{6e^2d_L^2}{\delta\mu^2}
                  +C\delta\mu^2\Lambda_\gamma^2\right)V
 \leq E^{\vac,(\delta\mu)}_{\Lambda,L}
 \leq E_\gamma(\delta\mu)\leq0.
 \label{eq:ct-vac}
\end{equation}
Thus the extensive estimate \eqref{eq:W} holds for this modified model
with the stated normal ordering. No additional infrared decoupling of
the original nonzero current couplings is imposed.
\end{theorem}
\begin{proof}
Conjugate by \(I\otimes U\). The fermionic terms are unchanged, the
photon Hamiltonian is given by \eqref{eq:proca-H}, and the interaction
contains the field \eqref{eq:proca-A}. Every current coefficient is a bounded
operator on the finite-dimensional fermionic space. A finite sum of such
coefficients times creation or annihilation operators is infinitesimally
bounded relative to the positive oscillator Hamiltonian. The Coulomb and
free fermionic terms are bounded for each fixed configuration and fixed
regulators. Kato--Rellich therefore gives self-adjointness on the stated
domain. Equivalently, the perturbation is infinitesimally form bounded,
so its form norm is equivalent to the oscillator form norm after adding a
constant. Finitely many positive-frequency oscillators have compact
form-domain embedding, proving compact resolvent. Charge conservation
reduces the operator to the indicated sectors.

Apply Lemma~\ref{lem:square} with frequencies \(\omega_k\) and coefficients
\[
 B_{k,\lambda}=
 \frac{e}{\sqrt{2V\omega_k}}\,
 \varepsilon_\lambda(k)\cdot\widehat j(k).
\]
Here and below \(|B|^2\) denotes \(B^*B\), not a commuting numerical square.
The result is
\begin{align*}
 \sum_{k,\lambda}\omega_k a^\dagger a
       +e\int j\cdot UAU^*
 &\geq-\frac{e^2}{2V}\sum_{k,\lambda}
       \frac{|\varepsilon_\lambda(k)\cdot\widehat j(k)|^2}{\omega_k^2}\\
 &\geq-\frac{e^2}{2\delta\mu^2}
             \int\sum_{i=1}^3j_i(x)^2\dd x\\
 &\geq-\frac{6e^2d_L^2V}{\delta\mu^2}.
\end{align*}
For the second step the transverse polarization matrix is a positive
contraction on the vector index, so its operator quadratic form is bounded
by the full vector square, regardless of commutation between currents.
Extending the Fourier sum and applying Parseval then gives the integral.
The last step is \eqref{eq:local-bounds}. Add \(H_D\geq0\), the
electrostatic estimate \eqref{eq:full-coulomb-bound}, and the scalar bound
in Lemma~\ref{lem:proca}. This proves \eqref{eq:ct-bound} and the lower
vacuum bound.

For the upper vacuum bound, use
\(\Omega_f\otimes U^*\Omega_\gamma\). Its free fermion and fully
Wick-ordered Coulomb expectations vanish at \(K=0\). Its current-field
expectation vanishes because the free-sea current is zero. Its photon
energy is exactly \(E_\gamma(\delta\mu)\).
\end{proof}

\begin{remark}[Soft-momentum matching and the remaining distinction]
\label{rem:ct-choice}
\begin{enumerate}[label=(\alph*)]
\item The deformed free-sea coefficient at quadratic order is
\[
q^2+\delta\mu^2-2e^2\chi_L(q;\varepsilon).
\]
Thus the choice
\(\delta\mu^2=\mu_{\rm ind}^2=2e^2\chi(0)\) cancels its constant part
in the infinite-volume, soft-momentum limit at order \(e^2\).
We call this \emph{static leading-order matching}. It is not an exact
finite-volume identity and does not prove a nonperturbative massless photon
or the Ward identities of a renormalized theory.
\item At this choice,
\[
 \frac{6e^2d_L^2}{\delta\mu^2}
 =\frac{3d_L^2}{\chi(0)}
 \longrightarrow\frac{\Lambda_f^3E_\Lambda}{2\pi^2}
 \quad(L\to\infty).
\]
The cancellation of \(\alpha\) in this particular bound is algebraic.
It is not a statement about cutoff-uniform renormalized parameters.
\item In the neutral-current trial construction, optimizing the displaced
massive oscillator coordinates replaces the massless photon gain by
\[
 -\frac{e^2J^2}{V(q^2+\delta\mu^2)},
 \qquad
 \frac{e^2J^2}{V(q^2+\delta\mu^2)}
 \leq\frac{e^2J^2}{V\delta\mu^2}=O(V).
\]
The scalar photon contribution is also \(O(V)\). The \(L^5\) mechanism
is therefore removed in this modified model, consistently with the theorem.
\item Theorem~\ref{thm:ct} needs only \(\delta\mu>0\). It gives no
uniformity as \(\delta\mu\downarrow0\) or as the ultraviolet cutoffs are
removed. It does not restore gauge covariance at nonzero momentum or
establish the physical massless renormalization conditions. The theorem
is an alternative extensive statement for a changed bare Hamiltonian;
it does not supersede the earlier infrared-cutoff result within its own
model. Comparing nuclear and interacting vacuum energies still requires
a relative estimate beyond \eqref{eq:ct-bound}--\eqref{eq:ct-vac}.
\end{enumerate}
\end{remark}

\subsection{An analytic small-box check for the retained self-image convention}

\begin{remark}[Small-box obstruction under both specified conventions]
\label{rem:madelung}
Let
\[
 r_{\rm self}=\lim_{x\to0}\bigl(G_1(x)-|x|^{-1}\bigr),\qquad
 \xi=-r_{\rm self}.
\]
We first prove \(\xi>0\) without relying on sampled Ewald values.
For \(\eta>0\), the Ewald identity for the zero-mean unit-torus Green
function is
\begin{align}
 G_1(x)
 &=\sum_{n\in\Z^3}\frac{\operatorname{erfc}(\eta|x+n|)}{|x+n|}
   +\frac1\pi\sum_{n\ne0}
       \frac{e^{-\pi^2|n|^2/\eta^2}}{|n|^2}
                  \cos(2\pi n\cdot x)-\frac\pi{\eta^2}.
 \label{eq:ewald-analytic}
\end{align}
Both sums converge absolutely away from lattice singularities. To verify
the normalization, write
\[
 \frac{\operatorname{erfc}(\eta|x|)}{|x|}
 =\frac2{\sqrt\pi}\int_\eta^\infty e^{-t^2|x|^2}\dd t.
\]
Its Fourier transform at \(k\ne0\) is
\(4\pi|k|^{-2}(1-e^{-|k|^2/(4\eta^2)})\), and its integral is
\(\pi/\eta^2\). Periodization and the reciprocal sum in
\eqref{eq:ewald-analytic} therefore give exactly the Fourier coefficients
\(4\pi/|2\pi n|^2\), with zero mean.

Set \(\eta=\sqrt\pi\), subtract \(|x|^{-1}\), and let \(x\to0\).
The singular real-space summand contributes \(-2\), while the last
constant in \eqref{eq:ewald-analytic} contributes \(-1\). Hence
\begin{equation}
 r_{\rm self}
 =\sum_{n\ne0}\frac{\operatorname{erfc}(\sqrt\pi|n|)}{|n|}
  +\frac1\pi\sum_{n\ne0}\frac{e^{-\pi|n|^2}}{|n|^2}-3.
 \label{eq:madelung-sign}
\end{equation}
The elementary Gaussian-tail bound
\(\operatorname{erfc}(t)\leq e^{-t^2}/(\sqrt\pi t)\), \(t>0\), yields
\begin{align*}
 r_{\rm self}
 &\leq\frac2\pi\sum_{n\ne0}\frac{e^{-\pi|n|^2}}{|n|^2}-3\\
 &\leq\frac2\pi
   \left[\left(1+2\sum_{j\geq1}e^{-\pi j^2}\right)^3-1\right]-3.
\end{align*}
Since \(\pi>3\) and \(e^3>16\), one has \(e^{-\pi}<1/16\), and
\[
 \sum_{j\geq1}e^{-\pi j^2}
 \leq\sum_{j\geq1}e^{-\pi j}<\frac1{15}.
\]
It follows, using only rational bounds, that
\[
 r_{\rm self}
 <\frac23\left[\left(\frac{17}{15}\right)^3-1\right]-3<0.
\]
This proves \(\xi>0\). It does not assert a sign or a maximum location for
the regular part throughout the unit cell.

If only the divergent Euclidean nuclear self-energy is removed, the
retained finite self-image term is
\(-\alpha\xi(2L)^{-1}\sum_iZ_i^2\). Now take exactly the small-box example
of \eqref{eq:small-box}: two nuclei with \(Z_1=Z_2=1\), only the fermion
momentum \(p=0\), and no retained nonzero photon modes. These charges are
compatible with the finite electron capacity. Zero mean and the positive
singularity of \(G_1\) imply that some \(r_*\ne0\) has \(G_1(r_*)<0\).
For \(R_1-R_2=Lr_*\), retaining the self-image term gives
\[
 E_{\Lambda,L}(R,Z)-E^{\vac}_{\Lambda,L}
 =2m+\frac\alpha L\bigl(G_1(r_*)-\xi\bigr)
 \longrightarrow-\infty\qquad(L\downarrow0).
\]
Under the convention removing the entire nuclear diagonal, the same
conclusion follows with \(\xi\) absent. Thus the lower restriction
\(L\geq L_0>0\) is necessary for an all-configurations relative theorem
under either of these two specified conventions. The photon deformation
does not change this example because the sufficiently small boxes retain
no photon modes.
\end{remark}
